\chapter{État de l'Art et Positionnement}

\section*{Introduction du chapitre}
Avant de concevoir une nouvelle solution, il est indispensable d'examiner les approches
existantes afin d'en comprendre les forces, les limites et les zones d'innovation possibles.
Ce chapitre dresse d'abord un panorama des principales échelles de triage normalisées
(CTAS, ESI, MTS, KTAS), puis des outils logiciels et des travaux de recherche mobilisant
l'apprentissage automatique pour le triage. Il se conclut par une analyse comparative
positionnant QuickTriage au regard de ces solutions.

\section{Analyse des solutions existantes}

\subsection{Les échelles de triage normalisées}
Afin de réduire la variabilité du triage, plusieurs échelles standardisées à cinq niveaux
ont été développées et largement adoptées à l'échelle internationale.

\begin{itemize}
  \item \textbf{CTAS} (\textit{Canadian Triage and Acuity Scale}) : échelle canadienne à
  cinq niveaux, associant à chaque niveau un délai cible de prise en charge ; elle s'appuie
  sur des plaintes types et des modificateurs (paramètres vitaux, douleur).
  \item \textbf{ESI} (\textit{Emergency Severity Index}) : échelle nord-américaine à cinq
  niveaux fondée sur un algorithme décisionnel qui combine la gravité immédiate et la
  consommation prévisible de ressources.
  \item \textbf{MTS} (\textit{Manchester Triage System}) : système britannique reposant sur
  des organigrammes (\textit{flowcharts}) par motif de consultation, aboutissant à un code
  couleur de priorité.
  \item \textbf{KTAS} (\textit{Korean Triage and Acuity Scale}) : adaptation coréenne de la
  CTAS, également à cinq niveaux, qui sert de référence à ce projet. Le niveau~1 correspond à
  une urgence vitale immédiate et le niveau~5 à un cas non urgent.
\end{itemize}

Le tableau~\ref{tab:echelles} synthétise les caractéristiques de ces échelles.

\begin{table}[h]
\centering
\renewcommand{\arraystretch}{1.3}
\begin{tabular}{p{2.2cm} p{2.2cm} p{2cm} p{6.5cm}}
\toprule
\textbf{Échelle} & \textbf{Origine} & \textbf{Niveaux} & \textbf{Principe de fonctionnement} \\
\midrule
CTAS & Canada & 5 & Plaintes types + modificateurs (constantes, douleur) \\
ESI  & États-Unis & 5 & Algorithme gravité + ressources prévisibles \\
MTS  & Royaume-Uni & 5 & Organigrammes par motif, code couleur \\
KTAS & Corée du Sud & 5 & Adaptation de la CTAS, niveau 1 = urgence vitale \\
\bottomrule
\end{tabular}
\caption{Comparaison des principales échelles de triage normalisées.}
\label{tab:echelles}
\end{table}

\subsection{Le regroupement métier adopté par QuickTriage}
Si l'échelle KTAS comporte cinq niveaux, l'usage opérationnel de triage privilégie souvent
un regroupement plus lisible. QuickTriage adopte un regroupement en \textbf{trois classes
métier}, présenté au tableau~\ref{tab:classes}, qui simplifie la lecture tout en préservant
l'essentiel de l'information clinique : la distinction entre l'urgence vitale, l'urgence
relative et le non-urgent.

\begin{table}[h]
\centering
\renewcommand{\arraystretch}{1.4}
\begin{tabular}{p{3.2cm} p{2.5cm} p{6.5cm}}
\toprule
\textbf{Classe métier} & \textbf{Niveaux KTAS} & \textbf{Délai de prise en charge} \\
\midrule
\textcolor{qtcritique}{\textbf{CRITIQUE}}     & KTAS 1--2 & Prise en charge immédiate ($<$ 15 min) \\
\textcolor{qturgent}{\textbf{URGENT}}         & KTAS 3    & Prise en charge rapide \\
\textcolor{qtnonurgent}{\textbf{NON\_URGENT}} & KTAS 4--5 & Prise en charge différée \\
\bottomrule
\end{tabular}
\caption{Regroupement des niveaux KTAS en trois classes métier dans QuickTriage.}
\label{tab:classes}
\end{table}

\subsection{Outils logiciels existants}
Sur le plan logiciel, l'informatisation du triage prend aujourd'hui plusieurs formes :
\begin{itemize}
  \item \textbf{Modules de triage intégrés aux DPI} (dossiers patients informatisés) : ils
  guident la saisie selon une échelle (souvent ESI ou MTS) mais reposent essentiellement sur
  des règles fixes et l'appréciation humaine, sans réelle prédiction statistique.
  \item \textbf{Calculateurs et applications de scores} : applications mobiles reproduisant
  les organigrammes d'une échelle donnée ; elles guident le soignant mais n'apprennent pas
  des données.
  \item \textbf{Solutions de recherche fondées sur l'IA} : prototypes académiques explorant
  des modèles prédictifs, rarement intégrés dans une chaîne applicative complète et
  déployable.
\end{itemize}

\subsection{Approches d'apprentissage automatique dans la littérature}
La littérature scientifique sur le triage assisté par apprentissage automatique a connu un
essor notable. Les travaux mobilisent typiquement des modèles tels que la régression
logistique, les forêts aléatoires, les machines à gradient boosté (dont XGBoost) ou encore
les réseaux de neurones, à partir des paramètres vitaux et de variables administratives. Les
principaux enseignements de ces travaux, qui ont guidé nos choix, sont les suivants :
\begin{itemize}
  \item Les méthodes d'\textbf{ensemble par boosting} (XGBoost et apparentés) offrent un
  excellent compromis entre performance et robustesse sur des données tabulaires
  hétérogènes.
  \item Le \textbf{déséquilibre des classes} (les cas critiques étant rares) constitue un
  obstacle majeur, justifiant des techniques de rééquilibrage telles que SMOTE.
  \item L'\textbf{évaluation} ne doit pas se limiter à l'exactitude globale ; le rappel sur
  les classes graves est l'indicateur de sécurité déterminant.
  \item La \textbf{fuite de données} est un piège fréquent qui conduit à surestimer
  artificiellement les performances ; sa détection est cruciale.
\end{itemize}

\section{Analyse comparative et positionnement de QuickTriage}

\subsection{Critères de comparaison}
Pour positionner QuickTriage, nous retenons les critères suivants : l'échelle de triage
utilisée, le degré d'automatisation, le recours effectif à l'apprentissage automatique,
l'intégration dans une application mobile, la prise en compte de la sécurité patient
(orientation vers le rappel des cas graves), et la complétude de la chaîne applicative
(de la donnée jusqu'au déploiement).

\subsection{Tableau comparatif}
Le tableau~\ref{tab:positionnement} confronte QuickTriage aux grandes familles de solutions
identifiées.

\begin{table}[h]
\centering
\renewcommand{\arraystretch}{1.4}
\small
\begin{tabular}{p{3.2cm} p{2.3cm} p{2.3cm} p{2.3cm} p{3.2cm}}
\toprule
\textbf{Critère} & \textbf{Échelles manuelles} & \textbf{Modules DPI} & \textbf{Protos IA (recherche)} & \textbf{QuickTriage} \\
\midrule
Échelle utilisée        & Variable & ESI/MTS & Variable & KTAS (3 classes) \\
Automatisation          & Aucune & Partielle (règles) & Modèle seul & De bout en bout \\
Apprentissage auto.     & Non & Non & Oui & Oui (XGBoost+Optuna) \\
Intégration mobile      & Non & Rare & Rare & Oui (Flutter) \\
Orientation sécurité    & Dépend de l'humain & Règles fixes & Variable & Rappel CRITIQUE optimisé \\
Chaîne complète         & --- & Partielle & Non & ML + J2EE + Mobile \\
\bottomrule
\end{tabular}
\caption{Positionnement comparatif de QuickTriage.}
\label{tab:positionnement}
\end{table}

\subsection{Synthèse du positionnement}
Il ressort de cette analyse que QuickTriage se distingue moins par l'emploi isolé d'une
technique que par la \textbf{cohérence et la complétude de sa chaîne}. Là où les outils
existants se cantonnent souvent à une seule dimension -- guidage par règles, ou modèle
prédictif sans intégration -- QuickTriage articule un pipeline d'apprentissage rigoureux
(détection de fuite, rééquilibrage, optimisation, calibration), une exposition
interopérable via ONNX et une API REST J2EE, ainsi qu'une application mobile pensée pour le
terrain. Surtout, le projet assume un parti pris clair de \textbf{sécurité patient} en
optimisant le rappel des cas critiques, ce qui le différencie des approches purement
orientées exactitude globale.

\section*{Conclusion du chapitre}
Cet état de l'art a permis de situer les échelles de triage de référence, de recenser les
outils existants et de dégager les enseignements de la littérature sur l'apprentissage
automatique appliqué au triage. Le positionnement comparatif a mis en évidence la valeur
ajoutée de QuickTriage : une solution intégrée, interopérable et orientée sécurité patient.
Forts de ce cadrage, nous pouvons désormais formaliser précisément les besoins du système,
objet du chapitre suivant.
